%% kapa: exam=1 round=36 year=2025 period=1 form=A subject=real_estate subject_ko=부동산학원론 seq=18 number=98
%% source: Q-Net 감정평가사 기출문제 (cst00309) artlSeq=5231525 file=raw/1008/5231525/2025년 제36회 감정평가사 1차시험 1교시 문제지 원본.pdf
%% conversion: kordoc 4.24.51 -> tools/gen1.py (auto); review=proofread-sonnet; flags=none
\begin{question}{98}
다음의 내용에 모두 해당하는 모기지(Mortgage)는?
\begin{framedtext}
○ 차입자가 금융기관에 지불하는 저당지불액이 증권 발행자를 통하여 투자자에게 그대로 전달되는 형태이다. \par
○ 기초자산인 모기지 풀(pool)의 현금흐름 및 저당권에 대한 소유권을 나타내는 지분형이다. \par
○ 금융기관이 유동화중개기관을 통해 발행할 수도 있고, 유동화중개기관을 통하지 않고 자체적으로 유동화전문회사(SPC)를 만들어 발행할 수도 있다. \par
○ 모기지 소유자는 채무불이행위험, 조기상환위험, 금리위험을 부담할 수 있다.
\end{framedtext}
\begin{choices}
  \item 저당이체증권(MPTS)
  \item 저당담보부채권(MBB)
  \item 지불이체채권(MPTB)
  \item 다계층증권(CMO)
  \item 상업용저당증권(CMBS)
\end{choices}
\end{question}
\answer{1}
